\documentclass[11pt,a4paper]{article}
\usepackage[T1]{fontenc}
\usepackage[utf8]{inputenc}
\usepackage{newtxtext}
\usepackage{amsmath,amsthm,mathtools}
\usepackage{newtxmath}
\usepackage[margin=27mm,headheight=15pt]{geometry}
\usepackage{microtype,booktabs,tabularx,array}
\usepackage[dvipsnames]{xcolor}
\usepackage{enumitem,fancyhdr,hyperref,bookmark}
\hypersetup{colorlinks=true,linkcolor=MidnightBlue,citecolor=MidnightBlue,
 urlcolor=MidnightBlue,pdftitle={Global Fueter Inversion: Period Obstructions, Topological Splitting, and Spherical Singularities},
 pdfauthor={OpenAI},pdfsubject={Research article extending the supplied quaternionic-analysis exposition}}
\urlstyle{same}
\numberwithin{equation}{section}
\newtheorem{theorem}{Theorem}[section]
\newtheorem{proposition}[theorem]{Proposition}
\newtheorem{lemma}[theorem]{Lemma}
\newtheorem{corollary}[theorem]{Corollary}
\theoremstyle{definition}
\newtheorem{definition}[theorem]{Definition}
\newtheorem{example}[theorem]{Example}
\theoremstyle{remark}
\newtheorem{remark}[theorem]{Remark}
\newcommand{\R}{\mathbb R}
\newcommand{\C}{\mathbb C}
\newcommand{\HH}{\mathbb H}
\newcommand{\HC}{\mathbb H_{\C}}
\newcommand{\SSS}{\mathbb S}
\newcommand{\SR}{\mathcal{SR}}
\newcommand{\AM}{\mathcal{AM}}
\newcommand{\cI}{\mathcal I}
\newcommand{\cD}{\mathcal D}
\newcommand{\lap}{\Delta_4}
\newcommand{\Per}{\operatorname{Per}}
\newcommand{\Res}{\operatorname{Res}}
\newcommand{\Rea}{\operatorname{Re}}
\newcommand{\Ima}{\operatorname{Im}}
\newcommand{\wind}{\operatorname{wind}}
\newcommand{\dd}{\,\mathrm d}
\newcommand{\ii}{\mathbf i}
\newcommand{\jj}{\mathbf j}
\newcommand{\kk}{\mathbf k}
\newcommand{\norm}[1]{\lVert #1\rVert}
\newcommand{\abs}[1]{\lvert #1\rvert}
\newcommand{\Kzero}{K_p^{(0)}}
\newcommand{\Kone}{K_p^{(1)}}
\newcommand{\Gzero}{G_p^{(0)}}
\newcommand{\Gone}{G_p^{(1)}}
\DeclareMathOperator{\dist}{dist}
\DeclareMathOperator{\supp}{supp}
\setlength{\emergencystretch}{2em}
\setlength{\parskip}{2pt}
\setlist{itemsep=2pt,topsep=4pt}
\pagestyle{fancy}
\fancyhf{}
\fancyhead[L]{\small Global Fueter inversion and spherical singularities}
\fancyhead[R]{\small\thepage}
\renewcommand{\headrulewidth}{0.3pt}
\title{\vspace{-10mm}\textbf{Global Fueter Inversion}\\[4pt]
\large Period obstructions, topological splitting,\\and spherical singularities}
\author{A research article developed from the supplied quaternionic-analysis manuscript}
\date{21 September 2026}
\begin{document}
\maketitle
\thispagestyle{empty}
\begin{abstract}
The local inverse of the quaternionic Fueter map is unique up to an affine
slice-regular function. We determine exactly how this ambiguity obstructs a
global inverse. For an axially monogenic function $g=P+IQ$, two explicit closed
quaternion-valued one-forms give necessary and sufficient period conditions for
$g=\Delta_4 f$. They also give the complete affine monodromy of every local
primitive. On finite-hole product domains the cokernel has right quaternionic
dimension twice the number of holes. On conjugation-symmetric domains meeting
the real axis, only pairs of exchanged holes contribute; holes fixed by
conjugation contribute nothing. Normalized logarithmic Fueter kernels split the
resulting exact sequence and yield quantitative non-approximation bounds and an
augmented rational approximation theorem.

The same invariants classify isolated nonreal spherical singularities. A
singularity of growth $O(\rho^{-N})$, where $\rho$ is distance to the singular
sphere and $N\geq1$, has a unique finite Laurent--logarithmic singular part of
real dimension $8N$. At the critical order $O(\rho^{-1})$, its two quaternionic
periods are exactly the coefficients of an affine surface-source density for
the Cauchy--Fueter operator. We compute the logarithmic energy divergence and
prove both the sharp $o(\rho^{-1})$ and the $L^2$ removability criteria. All
classification statements are proved in the axial class. Local Fueter inversion
and the basic logarithmic kernels are established mathematics; the proposed
contribution is the explicit global obstruction and singularity package.
\end{abstract}
\medskip
\noindent\textbf{Keywords.} Quaternionic analysis; inverse Fueter map; periods;
affine monodromy; axial monogenicity; spherical singularities; removable
singularities; approximation.

\medskip
\noindent\textbf{Status.} This is a research manuscript with complete analytic
proofs and supplementary formula checks, not a peer-reviewed publication or a
formal proof certificate. A targeted literature check did not locate the full
period, splitting, and singularity classification in the form proved here.
This is a qualified novelty assessment, not a guarantee of priority.
\newpage
\begingroup
\small
\setlength{\parskip}{0pt}
\tableofcontents
\endgroup
\newpage

\section{The question and the contribution}
\label{sec:introduction}

The supplied manuscript \cite{Source} ends its constructive inverse Fueter
theorem with a precise limitation: two closed one-forms can be integrated on a
simply connected planar set, but nonzero periods may prevent a global
primitive. The missing issue is therefore not local solvability. It is the
identification, realization, and analytic significance of the global
obstruction.

The inverse Fueter theory of Colombo, Sabadini, and Sommen \cite{CSS}, and its
integral development by Colombo, Pe\~na Pe\~na, Sabadini, and Sommen
\cite{CPSS}, already supplies local inverses and their polynomial ambiguity.
Perotti's Theorem~2 in \cite{Perotti} proves surjectivity when every component of
the underlying symmetric planar set is simply connected. We retain the same
notions of slice regularity and axial monogenicity, but remove that sufficient
topological hypothesis by replacing it with an exact criterion. The distinction
matters: some multiply connected domains are still globally surjective, while
others have a nonzero finite-dimensional cokernel.

\subsection{A useful change of unknowns}
Write a local slice-regular primitive as $f(x+Iy)=A(x,y)+IB(x,y)$, with $y>0$.
The conventional first integration constructs $C=B/y$. The conventional second
integration constructs $A$ using a form that already contains $C$. This makes
the second obstruction look dependent on a first choice of primitive. Instead
we use
\begin{equation}\label{eq:change-unknowns}
 C=\frac{B}{y},\qquad H=A-xC,\qquad f(x+Iy)=H+(x+Iy)C.
\end{equation}
Both $dC$ and $dH$ can then be written directly in terms of $g=\lap f$.
There are two independent periods, with no iterated integration or choice of
branch in their definition. The first is the coefficient of $q$ in the
monodromy; the second is its constant coefficient.

\subsection{What is proved}
The main results form one chain of consequences. The global criterion is
Theorem~\ref{thm:global}, including its real-axis version
Theorem~\ref{thm:real-axis}. The complete finite-hole obstruction space and its
explicit splitting are Theorem~\ref{thm:split}. A holomorphic detector converts
the periods into ordinary complex residues in
Theorem~\ref{thm:detector}. Theorem~\ref{thm:laurent-log} gives the finite-growth
classification at a nonreal sphere. Theorems~\ref{thm:source} and
\ref{thm:energy} identify the critical singularities as surface charges and
compute their energy. Quantitative approximation and removability consequences
are included because they show that the obstruction is analytically effective,
not merely an abstract dimension count.

These conclusions answer the extension question left open in the supplied
exposition. We do not identify that question with a named longstanding
conjecture in the published literature. The alternative requested here is to
prove substantial candidate-new theorems, with their relationship to established
work made explicit.

\subsection{What is not claimed as new}
Fueter's mapping theorem, its affine kernel in quaternionic dimension one,
the Vekua system, local inverse formulas, and the use of complex stems are not
new. In particular, the spherical integral kernels with primitives proportional
to $\arctan z$ and $z\arctan z$ already appear in \cite[Example~2]{CPSS}, with
attribution there to \cite{CSS}. Our logarithmic kernels are normalized and
translated versions of these established constructions, not newly discovered
special functions. What is proposed as the contribution is their role as a
complete period basis, together with the symmetry-sensitive global theorem and
the finite-order singularity, source, and energy classifications.

\section{Conventions and the differential system}
\label{sec:preliminaries}

\subsection{Quaternions, stems, and domains}
Let $\HH$ be the real quaternion algebra, with basis $1,\ii,\jj,\kk$ and
$\ii^2=\jj^2=\kk^2=\ii\jj\kk=-1$. Write
\[
 \SSS=\{I\in\HH:I^2=-1\},\qquad q=x+Iy,
 \quad x\in\R,\ y>0,\ I\in\SSS.
\]
We reserve $\iota$ for the central imaginary unit in
$\HC=\HH\otimes_{\R}\C$. Thus $\iota$ commutes with quaternionic coefficients
and is not a selected member of $\SSS$. The conjugation
$\kappa(A+\iota B)=A-\iota B$ acts only on this central complex factor.

For a conjugation-invariant planar open set $D$, put
\[
 \Omega_D=\{x+Iy:x+\iota y\in D,\ I\in\SSS\}.
\]
We consider connected circularizations of two kinds. A \emph{product domain}
comes from $D=U\cup\bar U$, where $U$ is a connected open subset of the upper
half-plane. A \emph{slice domain} comes from a connected symmetric $D$ meeting
$\R$. On a product domain, specifying a function on $U$ and reflecting it
specifies its stem on all of $D$.

A holomorphic stem is a holomorphic $F:D\to\HC$ satisfying
$F(\bar z)=\kappa F(z)$. Write $F=A+\iota B$, where $A,B$ are quaternion
valued. Its induced left slice-regular function is
\begin{equation}\label{eq:stem}
 \cI(F)(x+Iy)=A(x,y)+IB(x,y),\qquad
 A_x=B_y,\quad A_y=-B_x.
\end{equation}
These are the conventions of the supplied exposition and of the stem approach
\cite{GP}. All constant quaternionic coefficients in this article are multiplied
on the right. In particular, $\cI(z a+b)(q)=qa+b$ for $a,b\in\HH$.

The norm on $\HC$ used in coefficient estimates is the Euclidean norm
\[
 \norm{A+\iota B}_{\HC}^2=|A|^2+|B|^2.
\]
Multiplication by a central complex number of modulus one preserves this norm.
It is not the quaternionic norm of a slice evaluation.

\subsection{Axial monogenicity}
Our Cauchy--Fueter operator has no factor $1/2$:
\begin{equation}\label{eq:Dirac}
 \cD=\partial_{x_0}+\ii\partial_{x_1}
       +\jj\partial_{x_2}+\kk\partial_{x_3},\qquad
 \overline{\cD}\cD=\lap=\sum_{\mu=0}^3\partial_{x_\mu}^2.
\end{equation}
An \emph{axially monogenic} function has the form
$g(x+Iy)=P(x,y)+IQ(x,y)$, with smooth quaternion-valued $P,Q$, and satisfies
$\cD g=0$. On a slice domain we require smooth stem coefficients on $D$, with
$P$ even and $Q$ odd in $y$; the induced function is smooth across the real
axis. This is the monogenic slice-function class, denoted by $\AM(\Omega_D)$.
The space of slice-regular functions is denoted by $\SR(\Omega_D)$.

\begin{lemma}[Axial identities]\label{lem:axial}
For $y>0$,
\begin{align}
 \cD(P+IQ)&=P_x-Q_y-\frac{2Q}{y}+I(Q_x+P_y),\label{eq:axial-D}\\
 \lap(A+IB)&=A_{xx}+A_{yy}+\frac{2A_y}{y}
 +I\left(B_{xx}+B_{yy}+\frac{2B_y}{y}-\frac{2B}{y^2}\right).
 \label{eq:axial-lap}
\end{align}
Consequently, axial monogenicity is equivalent to
\begin{equation}\label{eq:Vekua}
 P_x-Q_y=\frac{2Q}{y},\qquad Q_x+P_y=0.
\end{equation}
For a slice-regular $f=A+IB$,
\begin{equation}\label{eq:Fueter-formula}
 \lap f=\frac{2A_y}{y}
       +I\left(\frac{2B_y}{y}-\frac{2B}{y^2}\right).
\end{equation}
The function $\lap f$ is axially monogenic.
\end{lemma}
\begin{proof}
Put $v=\Ima q$ and $I=v/y$. For $\ell=1,2,3$,
\[
 \partial_{x_\ell}y=\frac{x_\ell}{y},\qquad
 \partial_{x_\ell}I=\frac{e_\ell}{y}-\frac{v x_\ell}{y^3},
 \quad (e_1,e_2,e_3)=(\ii,\jj,\kk).
\]
It follows that
$\sum e_\ell\partial_{x_\ell}I=-2/y$,
$\Delta_{\R^3}I=-2I/y^2$, and
$\sum(\partial_{x_\ell}I)(\partial_{x_\ell}y)=0$.
The product rule gives \eqref{eq:axial-D} and \eqref{eq:axial-lap}.
Evaluating at $I$ and $-I$ separates the two coefficients, proving
\eqref{eq:Vekua}. The Cauchy--Riemann equations make both $A,B$ planar
harmonic, giving \eqref{eq:Fueter-formula}. If its two coefficients are $P,Q$,
direct differentiation using \eqref{eq:stem} gives \eqref{eq:Vekua}.
On a slice domain the identities extend across the real axis by smoothness.
\end{proof}
This calculation fixes all signs and normalizations needed below. For example,
$\lap q^2=-4$. At a real point, a holomorphic symmetric stem satisfies
\begin{equation}\label{eq:real-trace}
 (\lap\cI(F))(x)=-2F''(x).
\end{equation}
Indeed, its Taylor expansion has quaternionic coefficients, and only its
quadratic term contributes to the Laplacian at the center.

\section{Two explicit periods and the global inverse}
\label{sec:periods}

For $g=P+IQ\in\AM(\Omega_U)$ on a product domain, define
\begin{align}
 \alpha_0(g)&=-\frac{P}{2}\dd x+\frac{Q}{2}\dd y,
 \label{eq:alpha0}\\
 \alpha_1(g)&=\frac{xP+yQ}{2}\dd x+
              \frac{yP-xQ}{2}\dd y.
 \label{eq:alpha1}
\end{align}
Quaternion-valued forms are integrated componentwise over real curves.
For an oriented piecewise smooth closed curve $\gamma\subset U$, set
\begin{equation}\label{eq:period-def}
 a_\gamma(g)=\oint_\gamma\alpha_0(g),\qquad
 b_\gamma(g)=\oint_\gamma\alpha_1(g).
\end{equation}

\begin{lemma}[Closedness]\label{lem:closed}
Both forms in \eqref{eq:alpha0}--\eqref{eq:alpha1} are closed.
\end{lemma}
\begin{proof}
The coefficient of $dx\wedge dy$ in $d\alpha_0$ is $(Q_x+P_y)/2$.
The corresponding coefficient for $d\alpha_1$ is
\[
 \frac12\{y(P_x-Q_y)-2Q-x(Q_x+P_y)\}.
\]
Both vanish by \eqref{eq:Vekua}.
\end{proof}

\begin{theorem}[Complete period criterion]\label{thm:global}
Let $U$ be any connected open subset of $\{y>0\}$, and let
$g\in\AM(\Omega_U)$. The following conditions are equivalent:
\begin{enumerate}[label=\textup{(\roman*)}]
\item There is a single-valued $f\in\SR(\Omega_U)$ with $\lap f=g$.
\item Both $\alpha_0(g)$ and $\alpha_1(g)$ are exact on $U$.
\item $a_\gamma(g)=b_\gamma(g)=0$ for every closed curve $\gamma\subset U$.
\end{enumerate}
When these conditions hold, choose primitives $C,H$ of $\alpha_0,\alpha_1$.
Every inverse is
\begin{equation}\label{eq:global-inverse}
 f(x+Iy)=H(x,y)+(x+Iy)C(x,y),
\end{equation}
up to addition of $qa+b$ with $a,b\in\HH$.
\end{theorem}
\begin{proof}
Suppose first that $f=A+IB$ is an inverse, and define $C,H$ by
\eqref{eq:change-unknowns}. From \eqref{eq:Fueter-formula} and
$B_x=-A_y$,
\[
 C_x=-\frac P2,\qquad C_y=\frac Q2.
\]
Also,
\begin{align*}
 H_x&=A_x-C-xC_x=B_y-C+\frac{xP}{2}
       =\frac{yQ+xP}{2},\\
 H_y&=A_y-xC_y=\frac{yP-xQ}{2}.
\end{align*}
Thus $dC=\alpha_0$ and $dH=\alpha_1$.

Conversely, take primitives $C,H$, put $A=H+xC$ and $B=yC$, and compute
\[
 A_x=C+\frac{yQ}{2}=B_y,\qquad
 A_y=\frac{yP}{2}=-B_x.
\]
The resulting stem is holomorphic on $U$. Reflection supplies a holomorphic
symmetric stem on $U\cup\bar U$. Formula \eqref{eq:Fueter-formula} gives
$\lap f=g$.

A closed one-form has a primitive if and only if all its closed-curve integrals
vanish: fix a base point, integrate along a path, and use the loop condition
to prove path independence. This works for each real component. Finally, two
pairs of primitives differ by constant quaternions; their induced functions
differ by $qa+b$. Conversely, every such affine addition has zero Laplacian.
\end{proof}

\begin{remark}[Why there are two independent conditions]
The second period is not just the period of the form used to recover $A$ after
choosing $C$. It is already the period of an explicit form of $g$. The identity
$H=A-xC$ removes that apparent dependence. In particular, changing the first
primitive cannot remove a nonzero $b_\gamma$.
\end{remark}

\begin{corollary}[Affine monodromy]\label{cor:monodromy}
A local inverse continues along every path in $U$. Continuing it around a
closed curve $\gamma$ changes its value by
\begin{equation}\label{eq:monodromy}
 f^\gamma(q)-f(q)=q\,a_\gamma(g)+b_\gamma(g).
\end{equation}
The map $\gamma\mapsto(a_\gamma,b_\gamma)$ factors through
$H_1(U;\mathbb Z)$ and is additive. Equivalently, on the universal cover the
inverse is a single-valued function with this affine deck transformation law.
\end{corollary}
\begin{proof}
Integrate the two closed forms along paths from a fixed point. Their changes
under a loop are precisely their periods. Substituting in
\eqref{eq:global-inverse} proves the formula. Closed-form integration is
invariant under homology, and integration over concatenated paths is additive.
\end{proof}

This monodromy is not a quaternionic fractional-linear action. It is translation
by an element of the fixed eight-real-dimensional space
$\{q\mapsto qa+b:a,b\in\HH\}$. Noncommutativity has no hidden role in the
composition law because the coefficients remain on the right.

\section{Reflection, the real axis, and topology}
\label{sec:reflection}

Let $D$ now be connected, symmetric under
$\tau(x,y)=(x,-y)$, and meeting $\R$. Define $\alpha_0,\alpha_1$ by the same
formulas on all of $D$. No division by $y$ occurs in those definitions.

\begin{theorem}[Global inversion across the real axis]\label{thm:real-axis}
For $g\in\AM(\Omega_D)$, the forms $\alpha_0(g),\alpha_1(g)$ are smooth,
closed, and $\tau$-invariant. A global slice-regular inverse exists if and only
if both have zero periods on $D$. The inverse is unique modulo $qa+b$.
The obstruction belongs to
\begin{equation}\label{eq:invariant-cohom}
 H^1_{\mathrm{dR}}(D;\HH)^+\oplus H^1_{\mathrm{dR}}(D;\HH)^+,
\end{equation}
where the superscript denotes the $+1$ eigenspace of $\tau^*$.
\end{theorem}
\begin{proof}
Smoothness follows from the smooth stem coefficients. Closedness holds off
$\R$ by Lemma~\ref{lem:closed} and on $\R$ by continuity. The parity of $P,Q$
and the sign change of $dy$ prove $\tau^*\alpha_j=\alpha_j$.

If the periods vanish, choose primitives and replace them, if necessary, by
\[
 C_e=\tfrac12(C+C\circ\tau),\qquad
 H_e=\tfrac12(H+H\circ\tau).
\]
They have the same differentials and are even in $y$. Put
$A=H_e+xC_e$, $B=yC_e$. These have the required parities and satisfy the
Cauchy--Riemann equations off the real axis, hence everywhere by continuity.
They therefore define a holomorphic symmetric stem and the desired inverse.

For necessity, the odd smooth coefficient $B$ of any slice-regular inverse
has a smooth even quotient $B/y$ across $y=0$; for instance,
$B(x,y)/y=\int_0^1B_y(x,ty)\dd t$ near the axis. The computations in the
proof of Theorem~\ref{thm:global} extend by continuity and make both forms exact.
The same computation gives the affine kernel. Finally, invariance of the
forms gives \eqref{eq:invariant-cohom}.
\end{proof}

\subsection{A precise finite-hole setting}
For the finite-dimensional results, a \emph{finite-hole planar domain} will mean
a simply connected planar domain $V$ from which finitely many pairwise disjoint
compact holes have been removed. Each hole is either a point or the closure of
a Jordan domain, is contained in $V$, and is disjoint from the other holes.
We assume the remaining domain is connected. The outer domain $V$ may be
unbounded. In the product case $V\subset\{y>0\}$. In the slice case $V$ and
the family of holes are conjugation invariant, and the resulting domain meets
$\R$. These hypotheses include disks and half-planes with punctures, ordinary
annuli, and the punctured plane examples below.

Choose counterclockwise loops $\gamma_j$ with winding number one around the
$j$th hole and zero around the other holes. Their classes form a basis of the
first homology. One way to see this is to join the holes to the outer boundary,
or to a cut going to infinity, by disjoint arcs. Cutting along these arcs gives
a simply connected planar region; restoring each arc adds exactly one loop.
Replacing a point hole by a sufficiently small disk does not change this
argument. Only this elementary winding description of topology is needed.

\begin{lemma}[Which holes contribute]\label{lem:holes}
Suppose a symmetric finite-hole domain has $s$ holes fixed by conjugation and
$h$ pairs of holes exchanged by conjugation. Then
\[
 \dim_{\R}H^1_{\mathrm{dR}}(D;\R)^+=h.
\]
An invariant closed form has period zero around each fixed hole, and its
counterclockwise periods around a conjugate pair are negatives of one another.
\end{lemma}
\begin{proof}
Reflection reverses planar orientation. On winding generators it sends the
class around a fixed hole to its negative, and sends the class around one
member of an exchanged pair to the negative class around the other member.
For an invariant form $\alpha$,
$\int_{\tau\gamma}\alpha=\int_\gamma\tau^*\alpha=\int_\gamma\alpha$.
This proves the asserted relations. Conversely, those are the only relations
on a winding basis, leaving one real coordinate per exchanged pair.
\end{proof}

\begin{corollary}[Multiply connected surjectivity]\label{cor:annulus}
If every hole of a symmetric finite-hole slice domain is fixed by conjugation,
then $\lap:\SR(\Omega_D)\to\AM(\Omega_D)$ is surjective. In particular this
holds for $D=\{r<|z|<R\}$, with $0<r<R$.
\end{corollary}
\begin{proof}
All periods of both invariant forms vanish by Lemma~\ref{lem:holes}; apply
Theorem~\ref{thm:real-axis}.
\end{proof}

Thus the first Betti number alone does not measure the obstruction. The
conjugation action on the first homology is essential. An annulus has a
nonzero first Betti number but no invariant period obstruction.

\section{Logarithmic kernels and the split exact sequence}
\label{sec:kernels}

\subsection{Single-valued Fueter images of multivalued stems}
Fix $p=u+\iota v$ with $v\ne0$. On local branches away from $p$, write
\begin{equation}\label{eq:logstem}
 h_p(z)=\frac{\log(z-p)}{2\pi\iota}.
\end{equation}
A counterclockwise circuit around $p$ changes $h_p$ by $1$ and $zh_p$ by $z$.
These changes induce affine slice functions, whose Laplacians vanish. It is
therefore meaningful on $\{y>0\}\setminus\{p\}$ to define the single-valued
axially monogenic functions
\begin{equation}\label{eq:K-definition}
 \Kzero=\lap\cI(h_p),\qquad \Kone=\lap\cI(zh_p).
\end{equation}
The notation means the common result from local branches; it does not assert
a global logarithm or a global primitive.

Put $X=x-u$, $Y=y-v$, $R_p^2=X^2+Y^2$, and $L_p=\log R_p$. For
$K_p^{(j)}=P_j+IQ_j$, direct formulas are
\begin{align}
 P_0&=\frac{X}{\pi yR_p^2},&
 Q_0&=-\frac{Y}{\pi yR_p^2}+\frac{L_p}{\pi y^2},
 \label{eq:K0}\\
 P_1&=\frac{xX+yY}{\pi yR_p^2}+\frac{L_p}{\pi y},&
 Q_1&=\frac{X}{\pi R_p^2}-\frac{xY}{\pi yR_p^2}
                +\frac{xL_p}{\pi y^2}.
 \label{eq:K1}
\end{align}
Every term is single valued. In particular, the arguments of the logarithms
have disappeared.

\begin{proposition}[Normalized periods]\label{prop:K-periods}
For every closed curve $\gamma$ in the upper half-plane avoiding $p$,
\begin{align}
 (a_\gamma(\Kzero),b_\gamma(\Kzero))
   &=(0,\wind(\gamma,p)),\label{eq:K-period0}\\
 (a_\gamma(\Kone),b_\gamma(\Kone))
   &=(\wind(\gamma,p),0).\label{eq:K-period1}
\end{align}
\end{proposition}
\begin{proof}
Take a local argument $\theta=\arg(z-p)$. For $h_p$ the stem coefficients are
$A=\theta/(2\pi)$ and $B=-L_p/(2\pi)$. Hence
\[
 C_0=-\frac{L_p}{2\pi y},\qquad
 H_0=\frac{\theta}{2\pi}+\frac{xL_p}{2\pi y}.
\]
For $zh_p$ the analogous functions are
\[
 C_1=\frac{\theta}{2\pi}-\frac{xL_p}{2\pi y},\qquad
 H_1=\frac{(x^2+y^2)L_p}{2\pi y}.
\]
Their differentials are $\alpha_0,\alpha_1$ by
Theorem~\ref{thm:global}. Along a loop $L_p$ returns to its initial value and
$\theta$ increases by $2\pi\wind(\gamma,p)$. The displayed periods follow.
Differentiating the same stems with \eqref{eq:Fueter-formula} also proves
\eqref{eq:K0}--\eqref{eq:K1}.
\end{proof}

\subsection{The symmetric kernels}
For $v>0$, define on the upper half-plane
\begin{equation}\label{eq:G-definition}
 G_p^{(j)}=K_p^{(j)}-K_{\bar p}^{(j)},\qquad j=0,1.
\end{equation}
They extend to smooth axially monogenic functions on
$\HH\setminus S_p$, where
\[
 S_p=u+v\SSS.
\]
To prove this at the real axis, use local branches of
\begin{equation}\label{eq:paired-log}
 \ell_p(z)=\frac{\log(z-p)-\log(z-\bar p)}{2\pi\iota},\qquad
 \ell_p'(z)=\frac{v}{\pi((z-u)^2+v^2)}.
\end{equation}
Near any real point the derivative is holomorphic and real on the real axis.
A local primitive with a real integration constant is therefore symmetric.
It agrees with a branch of \eqref{eq:paired-log} up to a real constant. Its
Fueter image, or that of $z\ell_p$, gives the asserted smooth extension.
On overlapping branches the difference is affine, so these extensions agree.

The counterclockwise periods around $p$ are $(0,1)$ and $(1,0)$; around
$\bar p$ they are their negatives. The real traces, useful for checking
normalizations, are
\begin{align}
 G_p^{(0)}(x)&=\frac{4v(x-u)}{\pi((x-u)^2+v^2)^2},\label{eq:G0-real}\\
 G_p^{(1)}(x)&=\frac{4v\{u(x-u)-v^2\}}
                    {\pi((x-u)^2+v^2)^2}.\label{eq:G1-real}
\end{align}
These follow at once from \eqref{eq:real-trace} and \eqref{eq:paired-log}.
For $p=\iota$, $\ell_p$ agrees locally, up to a real constant, with
$\pi^{-1}\arctan z$. Thus these are the familiar arctangent-based inverse
Fueter kernels, with a normalization chosen to make their periods equal to
one; compare \cite[Example~2]{CPSS}.

\subsection{The complete obstruction space}
In a finite-hole product domain choose a point $p_j$ in the $j$th omitted hole
and take $\mathcal K_j^{(r)}=K_{p_j}^{(r)}$. In a symmetric finite-hole slice
domain choose one upper hole from each exchanged pair, a point $p_j$ in it,
and take $\mathcal K_j^{(r)}=G_{p_j}^{(r)}$. Let $h$ be the number of chosen
holes. In both cases let $\gamma_j$ be the corresponding positive winding
loops and define
\begin{equation}\label{eq:Per-finite}
 \Per(g)=\bigl((a_j,b_j)\bigr)_{j=1}^h,
 \quad a_j=a_{\gamma_j}(g),\quad b_j=b_{\gamma_j}(g).
\end{equation}

\begin{theorem}[Split global obstruction sequence]\label{thm:split}
For either finite-hole setting just described, the following sequence of
right quaternionic vector spaces is exact:
\begin{equation}\label{eq:exact-sequence}
 0\longrightarrow\HH^2
 \xrightarrow{\ (a,b)\mapsto qa+b\ }\SR(\Omega_D)
 \xrightarrow{\ \lap\ }\AM(\Omega_D)
 \xrightarrow{\ \Per\ }(\HH^2)^h\longrightarrow0.
\end{equation}
The final arrow has the explicit right-linear section
\begin{equation}\label{eq:section-J}
 J\bigl((a_j,b_j)_{j=1}^h\bigr)
   =\sum_{j=1}^h\bigl(\mathcal K_j^{(1)}a_j+
                       \mathcal K_j^{(0)}b_j\bigr).
\end{equation}
Consequently every $g\in\AM(\Omega_D)$ has a decomposition
\begin{equation}\label{eq:normal-form}
 g=\lap f+\sum_{j=1}^h
      \bigl(\mathcal K_j^{(1)}a_j+\mathcal K_j^{(0)}b_j\bigr),
\end{equation}
with uniquely determined period coefficients and with $f$ unique modulo
$qa+b$. The cokernel of $\lap$ has right quaternionic dimension $2h$, or real
dimension $8h$.
\end{theorem}
\begin{proof}
The affine-kernel statement and exactness at $\AM$ are the global period
criterion. In the slice case Lemma~\ref{lem:holes} says that the selected upper
periods are all the independent periods. Proposition~\ref{prop:K-periods} and
its paired version give
\[
 \Per\bigl(\mathcal K_j^{(1)}a+\mathcal K_j^{(0)}b\bigr)
   =(0,\ldots,0,(a,b),0,\ldots,0).
\]
Thus $\Per J$ is the identity. Subtract $J\Per(g)$ from $g$ and apply the
period criterion to the resulting zero-period function. Applying $\Per$ to
any putative decomposition proves uniqueness of its coefficients. The
dimension statement follows from the explicit isomorphism of the quotient
with $(\HH^2)^h$.
\end{proof}

\begin{remark}[Canonical quotient, noncanonical representatives]
Moving the point $p_j$ inside its omitted hole changes the kernel representative
but not its period vector. The difference is therefore the Fueter image of a
global slice-regular function. The quotient and its period coordinates do not
depend on those auxiliary pole locations; the section $J$ does.
\end{remark}

\begin{example}[A sphere removed from quaternionic space]\label{ex:removed-sphere}
Let $D=\C\setminus\{\iota,-\iota\}$, so
$\Omega_D=\HH\setminus\SSS$. There is one exchanged pair of holes, hence the
cokernel has real dimension eight. The functions $G_{\iota}^{(0)}$ and
$G_{\iota}^{(1)}$ are globally axially monogenic, but neither is the Laplacian
of a global slice-regular function on this domain. Their local primitives are
$\pi^{-1}\arctan z$ and $\pi^{-1}z\arctan z$; going once around $\iota$
adds $1$ or $z$. The local primitives are not global despite the global
single-valuedness of their Fueter images.
\end{example}

\section{A holomorphic detector and residue formulas}
\label{sec:detector}

The closed forms are directly usable with real coordinates. A second
formulation is convenient when meromorphic complex formulas are available.

\begin{theorem}[Second-derivative detector]\label{thm:detector}
For $g=P+IQ\in\AM(\Omega_U)$ define the $\HC$-valued function
\begin{equation}\label{eq:W}
 W_g(z)=-\frac12\bigl(P+yP_y+\iota yP_x\bigr).
\end{equation}
It is globally holomorphic. If $F$ is the stem of any local Fueter primitive,
then $W_g=F''$. For every closed curve $\gamma$,
\begin{equation}\label{eq:W-periods}
 \oint_\gamma W_g(z)\dd z=a_\gamma(g),\qquad
 \oint_\gamma zW_g(z)\dd z=-b_\gamma(g).
\end{equation}
On a slice domain $W_g$ extends holomorphically across $\R$ and is symmetric
under $\kappa$.
\end{theorem}
\begin{proof}
A local primitive exists on every small disk in the upper half-plane. In the
notation of Theorem~\ref{thm:global}, its complex derivative is
\[
 F'=A_x+\iota B_x=C+\frac{yQ}{2}-\iota\frac{yP}{2}.
\]
Differentiate with respect to $x$ and use $C_x=-P/2$ and $Q_x=-P_y$.
This gives \eqref{eq:W}, proving both holomorphy and independence of the local
primitive. Along an analytically continued branch, the monodromy of $F$ is
$za_\gamma+b_\gamma$, and that of $F'$ is $a_\gamma$. Integrating
$dF'=W_g\,dz$ gives the first formula. Integrating
$d(zF'-F)=zW_g\,dz$ gives the second. Near a real point the local construction
in Theorem~\ref{thm:real-axis} gives a holomorphic symmetric $F$, which proves
the final assertion.
\end{proof}

\begin{corollary}[Residue extraction at a nonreal puncture]\label{cor:residues}
Suppose $W_g$ is meromorphic at $p$ and has Laurent expansion
$W_g(z)=\sum_n c_n(z-p)^n$. For a small counterclockwise circle around $p$,
\begin{equation}\label{eq:residue-periods}
 a=2\pi\iota\,c_{-1},\qquad
 b=-2\pi\iota\,(p c_{-1}+c_{-2}).
\end{equation}
Although computed in $\HC$, both expressions belong to the real subspace
$\HH\subset\HC$. The two periods vanish exactly when
$c_{-1}=c_{-2}=0$.
\end{corollary}
\begin{proof}
Apply the componentwise complex residue theorem to
\eqref{eq:W-periods}. Membership in $\HH$ follows independently from the
real quaternion-valued integrals \eqref{eq:period-def}. Since
$2\pi\iota$ is invertible, the vanishing criterion follows.
\end{proof}

For the two logarithmic generators the detector is particularly simple:
\begin{equation}\label{eq:W-logs}
 W_{\Kone a+\Kzero b}(z)=\frac{1}{2\pi\iota}
 \left(\frac{a}{z-p}-\frac{pa+b}{(z-p)^2}\right).
\end{equation}
This follows by differentiating
$(za+b)\log(z-p)/(2\pi\iota)$ twice. Thus the only meromorphic terms that
obstruct a single-valued second primitive are precisely those that encode
Fueter monodromy.

\begin{remark}[The detector is not injective]
It would be incorrect to infer $g=0$ from $W_g=0$ on a product domain. An
affine holomorphic stem with arbitrary $\HC$ coefficients need not induce a
Fueter-harmonic slice function. For example, $F(z)=\iota$ gives
$f(x+Iy)=I$ and $\lap f=-2I/y^2$, although $F''=0$. The affine kernel of
$\lap$ consists only of $za+b$ with $a,b\in\HH$, not all complexified affine
stems. The proof of the detector theorem uses actual local primitives, not
an injectivity assertion.
\end{remark}

\section{Quantitative inversion and approximation}
\label{sec:approximation}

\subsection{Period bounds and a quantitative obstruction}
For a compact planar curve $\gamma$, let $\Omega_\gamma$ be its circularization,
and put
\[
 L_\gamma=\int_\gamma ds,\qquad
 R_\gamma=\int_\gamma |z|\dd s.
\]
These definitions also apply to curves crossing the real axis in a slice
domain.

\begin{proposition}[Norm bounds for the periods]\label{prop:bounds}
Let $M=\sup_{q\in\Omega_\gamma}|g(q)|$. Then
\begin{equation}\label{eq:period-bounds}
 |a_\gamma(g)|\le\frac{L_\gamma M}{2},\qquad
 |b_\gamma(g)|\le\frac{R_\gamma M}{2}.
\end{equation}
In particular, for every global single-valued slice-regular $f$,
\begin{equation}\label{eq:distance-bound}
 \sup_{\Omega_\gamma}|g-\lap f|\ \geq\
 \max\left\{\frac{2|a_\gamma(g)|}{L_\gamma},
             \frac{2|b_\gamma(g)|}{R_\gamma}\right\}.
\end{equation}
Terms with a zero denominator may be omitted; for a nonconstant closed curve
both displayed denominators are positive.
\end{proposition}
\begin{proof}
The parallelogram identity gives
\[
 |P|^2+|Q|^2=\tfrac12\bigl(|P+IQ|^2+|P-IQ|^2\bigr)\leq M^2.
\]
For a real unit tangent $(\dot x,\dot y)$, the norm of
$-P\dot x+Q\dot y$ is at most $(|P|^2+|Q|^2)^{1/2}$ by the Euclidean
Cauchy--Schwarz inequality. The two real coefficients of $P,Q$ in
$2\alpha_1(\dot\gamma)$ are
$x\dot x+y\dot y$ and $y\dot x-x\dot y$; their squared sum is $|z|^2$.
Integrating proves \eqref{eq:period-bounds}. Apply the same estimates to
$g-\lap f$, whose periods are those of $g$, to obtain
\eqref{eq:distance-bound}.
\end{proof}

The bound proves more than the failure of exact inversion: a nonzero period
prevents arbitrarily accurate uniform approximation on a compact circular
curve by Fueter images of global slice functions.

\subsection{A normalized right inverse on the zero-period space}
Fix $z_*=x_*+\iota y_*$ in $U$ and normalize $C(z_*)=H(z_*)=0$. For zero-period
$g$, the path-independent formula is
\begin{equation}\label{eq:path-inverse}
 f(q)=q\int_{z_*}^{z}\alpha_0(g)+\int_{z_*}^{z}\alpha_1(g),
 \qquad q=x+Iy,
\end{equation}
where the factors $q$ remain on the left of the quaternion-valued integrals.
For any chosen path $\sigma$ from $z_*$ to $z$, if
$M_\sigma=\sup_{\Omega_\sigma}|g|$, then
\begin{equation}\label{eq:path-bound}
 |f(q)|\leq\frac{M_\sigma}{2}
  \left(|q|L_\sigma+\int_\sigma|w|\dd s\right).
\end{equation}
This follows from the proof of Proposition~\ref{prop:bounds} applied to an
open path. The normalization means that $f$ vanishes on the full sphere over
$z_*$. It removes precisely the eight real affine degrees of freedom. On a
slice domain one may instead choose $z_*=x_*\in\R$; the same normalization
is $f(x_*)=\partial_c f(x_*)=0$.

\begin{proposition}[Closed range and continuous splitting]\label{prop:closed-range}
Within $\AM(\Omega_D)$ equipped with locally uniform convergence,
$\lap\SR(\Omega_D)$ is closed. In the finite-hole setting,
\[
 g\longmapsto g-J\Per(g)
\]
is a continuous projection onto this range. A base-point-normalized inverse
on the range is continuous for locally uniform convergence.
\end{proposition}
\begin{proof}
Every period is a continuous functional of $g$ on a compact circularized loop,
by \eqref{eq:period-bounds}. The range is the intersection of their kernels,
so is closed. In finite connectivity, $J$ is a finite sum of fixed smooth
functions times period coordinates and is continuous; $\Per J$ is the identity.

For the normalized inverse, cover a compact set by finitely many small disks
and choose finitely many connecting paths from the base point. All required
path integrals can then be taken in one larger compact set, with uniformly
bounded lengths. The estimate \eqref{eq:path-bound} applied to differences
proves locally uniform convergence of the inverses. The same construction,
with symmetric paths or even primitives, handles the real axis. Holomorphic
Cauchy estimates then also give convergence of derivatives on smaller compact
sets.
\end{proof}

\subsection{An augmented rational approximation theorem}
By a \emph{rational stem} we mean an $\HC$-valued rational function $R$ of the
central complex variable, with no poles in $D$ and with
$R(\bar z)=\kappa R(z)$. Its induced slice function $\cI(R)$ is globally
single valued and slice regular on $\Omega_D$.

\begin{theorem}[Approximation with the minimal obstruction space]
\label{thm:Runge}
In the finite-hole setting, let $(a_j,b_j)=\Per(g)$. There are rational stems
$R_n$, with poles outside $D$, such that
\begin{equation}\label{eq:augmented-Runge}
 g=\lim_{n\to\infty}\left[
   \lap\cI(R_n)+\sum_{j=1}^h
   \bigl(\mathcal K_j^{(1)}a_j+\mathcal K_j^{(0)}b_j\bigr)\right]
\end{equation}
locally uniformly, with all derivatives on compact subsets. Conversely, $g$
is a locally uniform limit of Fueter images of global slice-regular functions
if and only if $\Per(g)=0$. Any fixed right quaternionic linear augmentation
that makes such approximation possible for every $g$ must have dimension at
least $2h$ modulo $\lap\SR$.
\end{theorem}
\begin{proof}
Subtract the displayed kernel sum and use Theorem~\ref{thm:split} to obtain
$\lap f$ for a global slice-regular $f=\cI(F)$. The ordinary scalar Runge
theorem, applied to the four complex coordinates of $F$ on an exhaustion of
$D$, gives rational approximants with poles outside $D$. Replacing an
approximant $R$ by
\[
 \frac12\bigl(R(z)+\kappa R(\bar z)\bigr)
\]
preserves rationality, moves no pole into $D$, enforces stem symmetry, and
preserves convergence on symmetric compact sets. Complex Cauchy estimates give
convergence of all holomorphic derivatives locally. The induced slice
functions then converge with all real derivatives; near the real axis this
can be read from their convergent symmetric Taylor expansions. Applying $\lap$
proves \eqref{eq:augmented-Runge}.

The converse follows either from Proposition~\ref{prop:closed-range} or by
passing to the limit in every period. Finally, the images of any augmentation
must span the quotient $(\HH^2)^h$ under $\Per$. Its dimension is $2h$, so
fewer right quaternionic directions cannot suffice.
\end{proof}

The ordinary Runge theorem is the only approximation theorem imported in this
argument. The specifically quaternionic conclusion is the exact finite
augmentation required by the Fueter map; it is not a claim that all monogenic
functions are approximable by this axial construction.

\section{An isolated nonreal sphere: Laurent--logarithmic classification}
\label{sec:singularities}

Fix $p=u+\iota v$ with $v>0$ and choose $0<\varepsilon<v$. Let
\begin{equation}\label{eq:tube}
 U_p^*=\{z:0<|z-p|<\varepsilon\},\qquad
 T_p^*=\Omega_{U_p^*},\qquad S_p=u+v\SSS.
\end{equation}
For $q=x+Iy\in T_p^*$,
\[
 \rho=|z-p|=\sqrt{(x-u)^2+(y-v)^2}=\dist(q,S_p).
\]
The equality follows by minimizing the Euclidean distance from $q$ to
$u+vJ$, $J\in\SSS$; the minimizer is $J=I$. Define
\begin{equation}\label{eq:M-rho}
 M_g(\rho)=\sup_{\substack{|z-p|=\rho\\ I\in\SSS}}|g(x+Iy)|.
\end{equation}
All extensions below are through the \emph{whole sphere}, not through one point
of it. A regular part means an axially monogenic function extending to the full
tube $\Omega_{\{|z-p|<\varepsilon\}}$ after decreasing $\varepsilon$ if needed.

\subsection{The unrestricted local normal form}
\begin{proposition}[Holomorphic part plus two logarithmic modes]
\label{prop:local-normal}
For $g\in\AM(T_p^*)$, let $(a,b)$ be its periods around $p$. There is a
single-valued holomorphic $F_0:U_p^*\to\HC$ such that
\begin{equation}\label{eq:unrestricted-normal}
 g=\lap\cI(F_0)+\Kone a+\Kzero b.
\end{equation}
The coefficients $a,b$ are unique, and $F_0$ is unique up to $zc+d$ with
$c,d\in\HH$.
\end{proposition}
\begin{proof}
The punctured disk has one winding generator. Subtracting the two kernels with
the indicated coefficients kills its periods. Apply Theorem~\ref{thm:global}.
Uniqueness follows from periods and the affine-kernel statement.
\end{proof}
No growth hypothesis is needed here. The holomorphic $F_0$ may have an
essential singularity. A finite-growth assumption makes its singular part
finite and gives a precise dimension count.

\begin{lemma}[From growth of $g$ to growth of its detector]
\label{lem:detector-growth}
If $M_g(\rho)=O(\rho^{-N})$ for an integer $N\geq1$, then
\[
 W_g(z)=O(|z-p|^{-N-1}).
\]
Consequently $W_g$ is meromorphic at $p$ with pole order at most $N+1$.
\end{lemma}
\begin{proof}
Every real component of $g$ is harmonic because
$\overline{\cD}\cD g=\lap g=0$. A ball centered at a point of distance $\rho$
from $S_p$, with radius $\rho/2$, stays away from the sphere. Throughout that
ball the distance to the sphere is between $\rho/2$ and $3\rho/2$.
The harmonic interior gradient estimate therefore gives
$|\nabla g|=O(\rho^{-N-1})$, uniformly in $I$.

For completeness, the estimate follows directly from the ball mean-value
formula. For a harmonic real function $h$ on $B(a,r)$, differentiation under
its translated volume mean and the divergence theorem yield
\[
 \partial_\mu h(a)=\frac1{|B_r|}\int_{\partial B_r}h(a+\xi)n_\mu(\xi)\dd S,
 \qquad |\partial_\mu h(a)|\leq\frac4r\sup_{B(a,r)}|h|.
\]
A fixed pair of opposite slices recovers $P,Q$ and their planar derivatives
from values and derivatives of $g$, with bounded constants. Since $y$ stays
bounded away from zero, formula \eqref{eq:W} gives the desired estimate.
The Laurent coefficient estimate for a holomorphic function with this growth
forces all coefficients below degree $-N-1$ to vanish, componentwise.
\end{proof}

\begin{theorem}[Finite Laurent--logarithmic classification]
\label{thm:laurent-log}
Let $N\geq1$ be an integer and let $g\in\AM(T_p^*)$ satisfy
$M_g(\rho)=O(\rho^{-N})$. Then
\begin{equation}\label{eq:finite-normal}
 \boxed{\displaystyle
 g=g_{\mathrm{reg}}+
   \lap\cI\left(\sum_{k=1}^{N-1}(z-p)^{-k}c_{-k}\right)
   +\Kone a+\Kzero b,}
\end{equation}
where $c_{-k}\in\HC$, $a,b\in\HH$, and $g_{\mathrm{reg}}$ extends axially
monogenically through $S_p$. The singular coefficients, and hence the regular
part, are unique. Conversely every expression of this form has the stated
growth. Modulo extendible germs, the space of these singularities has real
dimension $8N$, or right quaternionic dimension $2N$.
\end{theorem}
\begin{proof}
Use Proposition~\ref{prop:local-normal} to subtract the logarithmic modes.
Each kernel in \eqref{eq:K0}--\eqref{eq:K1} is $O(\rho^{-1})$ because $y$
is bounded away from zero and $\log\rho=O(\rho^{-1})$. The remaining
zero-period function still has $O(\rho^{-N})$ growth. Write it as
$\lap\cI(F_0)$, where $F_0$ is single-valued holomorphic on the punctured disk.
Lemma~\ref{lem:detector-growth} says that $F_0''$ has a pole of order at most
$N+1$.

Expand $F_0(z)=\sum_{n\in\mathbb Z}d_n(z-p)^n$. For $n<0$, its second
derivative has coefficient $n(n-1)d_n$ at degree $n-2$. Since $n(n-1)\ne0$,
the pole-order bound implies $d_n=0$ for $n<-(N-1)$. Thus
\[
 F_0(z)=\sum_{k=1}^{N-1}(z-p)^{-k}c_{-k}+F_{\mathrm{reg}}(z),
\]
where $F_{\mathrm{reg}}$ is holomorphic on the full disk. Its induced Fueter
image is the required regular part. This argument also explains why terms
$\log(z-p)$ and $z\log(z-p)$ must be separated before applying an ordinary
Laurent expansion.

For uniqueness, subtract two decompositions. The period criterion makes the
two logarithmic coefficients zero. If the Fueter image of a finite principal
part extends monogenically, take a slice-regular primitive of that extension
on the full disk by the local inverse theorem. Subtract it from the principal
part. The difference has zero Fueter image on the punctured disk, so is
$zc+d$ with $c,d\in\HH$. An affine function cannot cancel a nonzero Laurent
principal part; all its coefficients therefore vanish.

For the converse, formula \eqref{eq:Fueter-formula} shows that the image of
$(z-p)^{-k}c$ is $O(\rho^{-k-1})$, and the kernels are
$O(\rho^{-1})$. Each $c_{-k}$ has eight real components, and $a,b$ have eight
real components together. Uniqueness proves that the resulting
$8(N-1)+8=8N$ parameters are independent modulo extendible germs.
\end{proof}

A convenient right quaternionic basis of the singular quotient consists of
$K_p^{(0)},K_p^{(1)}$ and, for $1\leq k\leq N-1$,
\begin{equation}\label{eq:sing-basis}
 \lap\cI((z-p)^{-k}),\qquad
 \lap\cI(\iota(z-p)^{-k}).
\end{equation}
The second generator at each order is essential: on a product domain,
complexified Laurent coefficients have two independent quaternionic parts.

\subsection{The sharp critical growth threshold}
\begin{corollary}[Critical order and subcritical removability]
\label{cor:critical}
If $M_g(\rho)=O(\rho^{-1})$, then
\begin{equation}\label{eq:critical-normal}
 g=g_{\mathrm{reg}}+\Kone a+\Kzero b.
\end{equation}
Such a function extends through $S_p$ if and only if $a=b=0$. In particular,
\begin{equation}\label{eq:small-o}
 M_g(\rho)=o(\rho^{-1})\quad\Longrightarrow\quad
 g\text{ extends axially monogenically through }S_p.
\end{equation}
The little-$o$ hypothesis cannot be replaced by $O(\rho^{-1})$.
\end{corollary}
\begin{proof}
The classification with $N=1$ has no Laurent principal part, proving the first
assertion. On a radius-$\rho$ circle,
$L_\gamma=2\pi\rho$ and
$R_\gamma\leq2\pi\rho(|p|+\rho)$. Proposition~\ref{prop:bounds} therefore
implies
\[
 |a|\leq\pi\rho M_g(\rho),\qquad
 |b|\leq\pi\rho(|p|+\rho)M_g(\rho).
\]
Under the little-$o$ hypothesis both periods vanish. A nonzero normalized
kernel has nonzero period and $O(\rho^{-1})$ growth, proving sharpness.
\end{proof}

The exponent here is one because the singular set is a nonreal two-sphere of
real codimension two. It is different from the isolated-point exponent three
in the Cauchy--Fueter removability theorem discussed in the supplied manuscript.
No statement for arbitrary nonaxial fields near arbitrary surfaces is being
inferred from this axial calculation.

\section{Spherical source densities and energy}
\label{sec:sources}

\subsection{The leading critical asymptotic}
Let $X=x-u$, $Y=y-v$ and $q_0(I)=u+Iv$. In the critical case
\eqref{eq:critical-normal}, the formulas for $K_p^{(0)},K_p^{(1)}$ give
\begin{equation}\label{eq:leading}
 g(x+Iy)=\frac{X-IY}{\pi v\rho^2}
                \bigl(q_0(I)a+b\bigr)+O(1+|\log\rho|).
\end{equation}
The remainder estimate is uniform in $I$ and in the polar angle of $(X,Y)$.
For example, the leading terms of $K_p^{(1)}$ are
\[
 P_1=\frac{uX+vY}{\pi v\rho^2}+O(1+|\log\rho|),\qquad
 Q_1=\frac{vX-uY}{\pi v\rho^2}+O(1+|\log\rho|).
\]
Together they are $(X-IY)(u+Iv)/(\pi v\rho^2)$. This verifies the order of
quaternionic multiplication in \eqref{eq:leading}.

Denote by $\delta_{S_p}$ the distribution given by Euclidean area measure on
$S_p$. A quaternion-valued density is placed on its left in the notation below;
thus $\sigma\delta_{S_p}$ acts on a real scalar test function $\varphi$ by
$\int_{S_p}\sigma(q)\varphi(q)\dd S_q$.

\begin{theorem}[The affine spherical charge]\label{thm:source}
An axially monogenic $g$ with $M_g(\rho)=O(\rho^{-1})$ is locally integrable
through $S_p$. Its distributional Cauchy--Fueter derivative in the full tube is
\begin{equation}\label{eq:source-density}
 \boxed{\displaystyle
 \cD g=\frac{2}{v}(q a+b)\,\delta_{S_p},}
\end{equation}
where $(a,b)$ are its two periods. In particular, the critical singular quotient
is identified with the space of affine quaternionic surface densities
$q\mapsto(2/v)(qa+b)$ on $S_p$.
\end{theorem}
\begin{proof}
In tubular coordinates
\[
 x=u+\rho\cos\theta,\qquad y=v+\rho\sin\theta,
 \qquad 0\leq\theta<2\pi,
\]
the volume measure is $y^2\rho\,d\rho\,d\theta\,d\sigma(I)$, where
$d\sigma$ is area measure on the unit sphere $\SSS$. Thus $O(\rho^{-1})$
is locally integrable.

Remove the tube $\rho<\delta$ and integrate the definition of the
distributional derivative by parts against a real scalar smooth compactly
supported $\varphi$. Since $\cD g=0$ off the sphere, the result is the limit of
the flux over the boundary of the removed tube, with normal directed away from
$S_p$. This normal and the area measure are
\[
 n_\delta=\cos\theta+I\sin\theta
         =\frac{X+IY}{\delta},\qquad
 dS=y^2\delta\,d\theta\,d\sigma(I).
\]
The sign is positive: the outward normal of the punctured domain on its inner
boundary is $-n_\delta$.

By \eqref{eq:leading},
\[
 n_\delta g=\frac{q_0(I)a+b}{\pi v\delta}
                 +O(1+|\log\delta|).
\]
After multiplication by $dS$, the remainder tends to zero in the flux integral.
Also $y^2\to v^2$ and $\varphi(x+Iy)\to\varphi(q_0(I))$ uniformly. The
limiting distribution is therefore
\[
 \int_{\SSS}2v\bigl(q_0(I)a+b\bigr)\varphi(q_0(I))\dd\sigma(I).
\]
Since $dS_{S_p}=v^2d\sigma(I)$, this is exactly
\eqref{eq:source-density}. The map from $(a,b)$ to the density is injective:
subtracting its values at $u+Iv$ and $u-Iv$ forces $a=0$, and then $b=0$.
Surjectivity onto the displayed affine density space follows from the kernels.
\end{proof}

\begin{corollary}[Total flux and first angular moment]\label{cor:moments}
For $\sigma(q)=(2/v)(qa+b)$ on $S_p$, define
\[
 M_0=\int_{S_p}\sigma(q)\dd S_q,\qquad
 M_1=\int_{S_p}(q-u)\sigma(q)\dd S_q.
\]
Then
\begin{equation}\label{eq:moments}
 M_0=8\pi v(ua+b),\qquad M_1=-8\pi v^3a.
\end{equation}
Thus the charge moments recover both periods:
\begin{equation}\label{eq:moment-inversion}
 a=-\frac{M_1}{8\pi v^3},\qquad
 b=\frac{M_0}{8\pi v}-ua.
\end{equation}
\end{corollary}
\begin{proof}
Use $q=u+Iv$, $|S_p|=4\pi v^2$, and $\int_{\SSS}I\dd\sigma=0$.
For the first moment, $(q-u)^2=-v^2$ and the remaining terms have zero spherical
average. The displayed factors follow. The weight in $M_1$ is multiplied on
the left, as required by the calculation.
\end{proof}

A total flux alone does not classify a spherical singularity. For $u=0$, $b=0$
and $a\ne0$, it vanishes although the singularity is nonremovable. The first
angular moment detects exactly the missing affine slope. This is a concrete
reason to retain both period invariants rather than introduce only one
quaternionic ``residue.''

\subsection{The exact logarithmic energy coefficient}
\begin{theorem}[Critical energy divergence]\label{thm:energy}
If $g$ has the critical decomposition \eqref{eq:critical-normal}, then, for a
fixed sufficiently small $\varepsilon_0>0$,
\begin{equation}\label{eq:energy}
 \int_{\delta<\dist(q,S_p)<\varepsilon_0}|g(q)|^2\dd V_q
 =8\bigl(|ua+b|^2+v^2|a|^2\bigr)
          \log\frac{\varepsilon_0}{\delta}+O(1)
 \quad(\delta\downarrow0).
\end{equation}
The coefficient is strictly positive for every nonremovable critical
singularity.
\end{theorem}
\begin{proof}
Write $X=\rho\cos\theta$, $Y=\rho\sin\theta$. Multiplicativity of the
quaternionic norm gives
\[
 \left|\frac{X-IY}{\pi v\rho^2}(q_0(I)a+b)\right|^2
     =\frac{|q_0(I)a+b|^2}{\pi^2v^2\rho^2}.
\]
The cross term with the remainder in \eqref{eq:leading} is
$O(\rho^{-1}(1+|\log\rho|))$, and the squared remainder is
$O((1+|\log\rho|)^2)$. Both are integrable against the tubular volume measure
near zero. Replacing $y^2$ by $v^2$ in the leading energy changes it by an
integrable term as well. Finally,
\[
 \int_{\SSS}|(u+Iv)a+b|^2\dd\sigma(I)
       =4\pi\bigl(|ua+b|^2+v^2|a|^2\bigr),
\]
because the term linear in $I$ averages to zero. Integration in $\theta$
contributes $2\pi$, and the remaining radial integral is
$\int_\delta^{\varepsilon_0}d\rho/\rho$. The total coefficient is
$(2\pi/\pi^2)4\pi=8$. Positivity follows from $v>0$.
\end{proof}

This energy formula is independent of the regular part. It assigns a positive
definite quadratic norm to the eight-real-dimensional critical obstruction
space, with a direct geometric meaning as the coefficient of logarithmic
energy concentration around the sphere.

\section{An integrability removability theorem}
\label{sec:L2}

The growth classification also yields an integrability statement without any
assumed pointwise growth bound.

\begin{lemma}[A nonzero Laurent mode cannot have finite $L^2$ energy]
\label{lem:pole-energy}
Suppose the finite normal form \eqref{eq:finite-normal} has a highest nonzero
Laurent coefficient $c_{-m}$, with $m\geq1$. Then $g$ is not locally in $L^2$
near $S_p$.
\end{lemma}
\begin{proof}
For a holomorphic stem $F=A+\iota B$, formula \eqref{eq:Fueter-formula} can
also be written
\begin{equation}\label{eq:first-derivative-leading}
 \lap\cI(F)=\frac{2I}{y}\cI(F')-\frac{2IB}{y^2}.
\end{equation}
For $F=(z-p)^{-m}c_{-m}$, the first term has order $\rho^{-m-1}$ and the
second has order at most $\rho^{-m}$. Replacing $y$ by $v$ in the first term
also changes it by only $O(\rho^{-m})$. All lower Laurent terms and the
logarithmic modes are of lower order than $\rho^{-m-1}$.

If $c=A_0+\iota B_0\ne0$, the spherical average satisfies
\[
 \frac1{4\pi}\int_{\SSS}|\cI(e^{\iota t}c)(I)|^2\dd\sigma(I)
   =\norm{c}_{\HC}^2.
\]
Here $\cI(e^{\iota t}c)(I)$ means evaluation of the constant complexified
coefficient at $I$; the identity follows by expanding its two quaternionic
parts and averaging the term linear in $I$. The central rotation preserves
their squared norm. Hence the squared leading term, integrated in $I$ and
$\theta$, is a strictly positive constant times $\rho^{-2m-2}$, while all
lower terms have smaller angular $L^2$ norm. The tubular volume contributes
$\rho\,d\rho$, so the energy dominates a positive multiple of
$\int_0 \rho^{-2m-1}\dd\rho$, which diverges.
\end{proof}

\begin{theorem}[$L^2$ removability at a nonreal sphere]\label{thm:L2}
Let $g\in\AM(T_p^*)$. If $g$ is locally square integrable through $S_p$, then
it extends axially monogenically through $S_p$. The exponent is sharp among
$L^p$ conditions: nonremovable critical kernels belong locally to $L^p$ for
every $0<p<2$.
\end{theorem}
\begin{proof}
Choose a fixed full tube on which the $L^2$ norm is finite. Harmonicity of the
components and the volume mean-value formula, followed by Cauchy--Schwarz,
give, on a ball of radius $\rho/2$ centered at $q$,
\[
 |g(q)|\leq |B_{\rho/2}|^{-1/2}\norm{g}_{L^2(B(q,\rho/2))}
       \leq C\rho^{-2}.
\]
For sufficiently small $\rho=\dist(q,S_p)$ these balls lie in the fixed
punctured tube, and the constant is uniform in $q$. Thus the classification
Theorem~\ref{thm:laurent-log} applies with $N=2$. There is at most one Laurent
coefficient $c_{-1}$, which must vanish by Lemma~\ref{lem:pole-energy}.
The critical energy formula then forces $a=b=0$. What remains is the
extendible regular part.

Conversely, a critical kernel is $O(\rho^{-1})$, so its $p$th power is
integrable against $\rho\,d\rho$ when $p<2$. Its nonzero period makes it
nonremovable. This proves sharpness.
\end{proof}

The proof separates two possible failures of removability: a meromorphic
principal part creates a power-law divergence, while nonzero affine monodromy
creates a logarithmic divergence. Finite $L^2$ energy excludes both.

\section{Worked examples and consequences}
\label{sec:examples}

\subsection{A pure constant monodromy}
Take $p=u+\iota v$ and $g=K_p^{(0)}b$, $b\ne0$, on a punctured tubular
neighborhood. Its period pair is $(0,b)$, its critical density is
$(2/v)b$ times area measure on $S_p$, and its energy coefficient is $8|b|^2$.
The local primitive $h_p(z)b$ returns with the constant increment $b$.
No Laurent principal part of a single-valued holomorphic stem can produce this
period, regardless of how high a pole is allowed.

\subsection{A pure slope invisible to total flux}
For $u=0$, let $g=K_p^{(1)}a$ with $a\ne0$. The period pair is $(a,0)$.
The density is $2Ia$ on $S_p=v\SSS$, whose integral is zero. Nevertheless,
\[
 M_1=-8\pi v^3a,\qquad
 \int_{\delta<\rho<\varepsilon_0}|g|^2\dd V
       =8v^2|a|^2\log(\varepsilon_0/\delta)+O(1).
\]
Thus zero total flux does not imply removability even at the weakest
nonremovable growth order. The obstruction is an affine monodromy, not a
higher-order pole hidden in the regular part.

\subsection{A pole with no period obstruction}
Let $F(z)=(z-p)^{-1}c$ with $c\in\HC\setminus\{0\}$, and set
$g=\lap\cI(F)$. It has zero periods because it already has a single-valued
slice-regular primitive. It is nevertheless singular on $S_p$, of growth
$O(\rho^{-2})$. Its detector is
\[
 W_g(z)=2(z-p)^{-3}c,
\]
which has no simple or double Laurent term and hence no period by
Corollary~\ref{cor:residues}. This distinguishes two notions:
\emph{global invertibility on a punctured domain} does not imply
\emph{removability through its missing sphere}.

\subsection{Several removed spheres}
Let $p_j=u_j+\iota v_j$ ($v_j>0$), $1\leq j\leq h$, be distinct points, and take
\[
 D=\C\setminus\{p_1,\bar p_1,\ldots,p_h,\bar p_h\}.
\]
Then $\Omega_D=\HH\setminus\bigcup_j S_{p_j}$, and every axially monogenic
function on this complement has a unique period part
\[
 \sum_{j=1}^h\bigl(G_{p_j}^{(1)}a_j+G_{p_j}^{(0)}b_j\bigr)
\]
modulo a global Fueter image. If it has at most critical growth near each
sphere, the local distributional source is
\[
 \cD g=\sum_{j=1}^h\frac{2}{v_j}(q a_j+b_j)\delta_{S_{p_j}}
\]
in any region where those are its only singularities. The periods can be measured on disjoint small loops, so charges at different
spheres are independent.

\subsection{Fixed and paired holes have different effects}
Start with a disk centered on the real axis. Removing a smaller concentric
closed disk creates one fixed hole but no Fueter obstruction. Removing instead
two small conjugate disks off the real axis creates one exchanged pair and an
eight-real-dimensional obstruction space. Removing both kinds simultaneously
leaves that dimension equal to eight. The geometric sizes of the holes affect
norm estimates and convenient kernel representatives, but not the quotient
dimension.

\section{Verification, scope, and further questions}
\label{sec:verification}

\subsection{What the supplementary computation checks}
The accompanying script \texttt{verify\_results.py} uses exact symbolic
simplification to check the Cauchy--Riemann equations, the displayed kernel
formulas, both Vekua equations, all four primitive derivatives, and the
holomorphic detectors for $K_p^{(0)}$ and $K_p^{(1)}$. Nine further symbolic checks use generic quaternionic right coefficients
to verify the charge moments and the energy quadratic form. There are 31
symbolic zero-residual checks in total. At 50-decimal-digit precision it also checks eight
period pairs: enclosing and nonenclosing loops for unpaired kernels and upper
and lower loops for paired kernels. Their maximum errors are below
$10^{-42}$. Four normal-circle flux calculations check convergence to the
surface density in Theorem~\ref{thm:source}; four further calculations check
the limiting logarithmic energy slope for noncommuting right coefficients.
The resulting version information,
values, and residuals are recorded in \texttt{verification\_results.json}.

These tests check arithmetic, signs, normalization, and numerical consistency.
They do not establish the global theorems, the topology argument, the Laurent
classification, or priority. Those claims rest on the proofs in this article.
No theorem prover was used, and no formal verification is claimed.

\subsection{Literature boundary and novelty assessment}
The targeted search was made on 21 September 2026. It included the established
inverse Fueter sources \cite{CSS,CPSS}, the global simply connected statement
in \cite{Perotti}, and searches for inverse Fueter periods, monodromy,
multiply connected inversion, and axial spherical singularities. Current
related preprints \cite{DS26,GGBS26} were also checked for scope. The first
concerns a polyanalytic variation of inverse Fueter theory and the second
operator connections; they are not the basis of any proof here.

No identical complete statement of the two-form obstruction, the
reflection-sensitive split sequence, or the $8N$ spherical singularity quotient
was located in that search. This supports treating those formulations and
their source/energy consequences as candidate original contributions, but does
not exclude an earlier treatment in the broader literature on Vekua systems,
Clifford analysis, or monodromy. In particular, standard cohomological
principles explain why a local inverse with an affine kernel should have
period obstructions. The value of the present result is the explicit complete
calculation, its realization by known kernels, and the further classifications,
not a claim to have invented that general principle.

\subsection{Boundaries of the proved results}
The global zero-period criterion is proved for arbitrary connected planar
bases of the two stated types. The explicit finite-dimensional surjectivity
onto all period vectors is proved for the finite-hole class specified in
Section~\ref{sec:reflection}; no infinite series of kernel generators is
silently assumed for infinitely connected domains. Rational approximation uses
ordinary Runge approximation and does not assert arbitrary prescribed boundary
values. All spherical singularity results concern axially monogenic functions
near a \emph{nonreal} sphere of positive radius. Real points and nonaxial fields
are different local problems.

Several extensions remain natural research questions. In infinitely connected
domains one can ask for a convergent kernel splitting with sharp norm control.
In higher odd spatial dimensions the local inverse kernel is a larger polynomial
space, suggesting higher polynomial monodromy and a corresponding family of
moment forms; that extension is not proved here. A further problem is a
weighted Hilbert-space theory in which the explicit critical energy norm
controls a canonical inverse after period subtraction. These are directions
suggested by the proved structure, not additional conclusions of this paper.

\section*{Conclusion}
\addcontentsline{toc}{section}{Conclusion}
The global Fueter inversion problem becomes explicit after replacing the two
primitive coefficients by $C=B/y$ and $H=A-xC$. Their differentials are closed
forms of the given monogenic function, and their periods give all affine
monodromy. This yields an exact global criterion, an explicit finite-hole
splitting, and quantitative restrictions on approximation. At a nonreal sphere,
the same two invariants are the logarithmic part of a finite Laurent--logarithmic
classification; at critical growth they are exactly an affine surface charge
and the coefficients of a positive logarithmic energy divergence. The local
inverse construction in the supplied exposition thereby extends to a global
and singularity-sensitive theory with fully specified topology and
multiplication conventions.

\appendix
\section{A compact audit of signs, dimensions, and normalizations}
\label{app:audit}
The central convention is $\cD=\partial_0+\ii\partial_1+\jj\partial_2+
\kk\partial_3$ and $\lap q^2=-4$. The affine monodromy is $qa+b$ with
coefficients on the right. Thus the slope period is $\oint\alpha_0$ and the
constant period is $\oint\alpha_1$; swapping them would reverse the roles of
the two logarithmic kernels.

For a local logarithmic primitive,
\[
 F(z)=\frac{(za+b)\log(z-p)}{2\pi\iota},\qquad
 F''(z)=\frac1{2\pi\iota}
   \left(\frac a{z-p}-\frac{pa+b}{(z-p)^2}\right).
\]
Hence $\oint F''dz=a$ and $\oint zF''dz=-b$, agreeing with the real-form
calculation. Reflection changes the sign of a counterclockwise winding loop,
which is why a fixed hole has no invariant period.

At a sphere $S_p=u+v\SSS$, a radius-$\rho$ normal circle has length $2\pi\rho$;
the sphere has area $4\pi v^2$; and the tube boundary element is
$y^2\rho\,d\theta\,d\sigma$. These give, respectively, the period estimates,
the source density $(2/v)(qa+b)$, and the energy coefficient
$8(|ua+b|^2+v^2|a|^2)$. No $2\pi^2$ appears in this codimension-two normal-circle
calculation. The $2\pi^2$ in the point Cauchy--Fueter kernel is the area of a
unit three-sphere and belongs to a different geometry.

Finally, at growth $O(\rho^{-N})$ there are $N-1$ complexified Laurent
coefficients, each of real dimension eight, plus two quaternionic logarithmic
coefficients, together of real dimension eight. The total is $8N$, not
$4N$, $8N+8$, or a count based only on isolated point residues.

\begin{thebibliography}{10}
\addcontentsline{toc}{section}{References}

\bibitem{Source}
Supplied manuscript,
\emph{Classical Complex Theorems over the Quaternions: Slice-regular and
Cauchy--Fueter analogues, with proofs}, 20 September 2026.
Unpublished expository source supplied as \texttt{quaternionic\_analysis.tex}.
The present article develops the global extension question following its
constructive axial inverse theorem.

\bibitem{GP}
R.~Ghiloni and A.~Perotti,
\emph{Slice regular functions on real alternative algebras},
Advances in Mathematics \textbf{226} (2011), 1662--1691.
\href{https://arxiv.org/abs/1008.4318}{arXiv:1008.4318}.

\bibitem{CSS}
F.~Colombo, I.~Sabadini, and F.~Sommen,
\emph{The inverse Fueter mapping theorem},
Communications on Pure and Applied Analysis \textbf{10} (2011), 1165--1181.
\href{https://doi.org/10.3934/cpaa.2011.10.1165}{doi:10.3934/cpaa.2011.10.1165}.

\bibitem{CPSS}
F.~Colombo, D.~Pe\~na Pe\~na, I.~Sabadini, and F.~Sommen,
\emph{A new integral formula for the inverse Fueter mapping theorem},
Journal of Mathematical Analysis and Applications \textbf{417} (2014), 112--122.
\href{https://doi.org/10.1016/j.jmaa.2014.03.016}{doi:10.1016/j.jmaa.2014.03.016};
\href{https://arxiv.org/abs/1302.0685}{arXiv:1302.0685}.

\bibitem{Perotti}
A.~Perotti,
\emph{A local Cauchy integral formula for slice-regular functions},
Computational Methods and Function Theory \textbf{24} (2024), 185--203.
Published online 30 May 2023.
\href{https://doi.org/10.1007/s40315-023-00485-5}{doi:10.1007/s40315-023-00485-5}.

\bibitem{DS26}
A.~De Martino and S.~Pinton,
\emph{A variation of the inverse Fueter theorem and the generalized
polyanalytic Cauchy--Kovalevskaya extension of order 2}, preprint (2026).
\href{https://arxiv.org/abs/2608.07711}{arXiv:2608.07711}.

\bibitem{GGBS26}
J.~O.~Gonz\'alez-Cervantes, D.~Gonz\'alez-Campos, J.~Bory-Reyes, and B.~Schneider,
\emph{A connection between quaternionic slice regular and Fueter
hyperholormorphic functions}, preprint (2026).
\href{https://arxiv.org/abs/2607.19496}{arXiv:2607.19496}.

\end{thebibliography}
\end{document}
